%% kapa: exam=1 round=30 year=2019 period=1 form=A subject=civil_law subject_ko=민법 seq=7 number=7
%% source: Q-Net 감정평가사 기출문제 (cst00309) artlSeq=5205205 file=extracted/1008/5205205/2019년 제30회 감정평가사 1차시험 문제지/2019년 제30회 감정평가사 1차시험 문제지/2019년도 제30회 감정평가사 1차 1교시 A형.hwp
%% conversion: kordoc 4.24.51 -> tools/gen1.py (auto); review=proofread-sonnet; flags=none
\begin{question}{7}
민법상 법인에 관한 설명으로 옳은 것을 모두 고른 것은? (다툼이 있으면 판례에 따름)
\begin{framedtext}
ㄱ. 재단법인의 설립을 위해 부동산의 출연이 행해진 경우, 그 부동산의 소유권은 그 출연 시에 곧바로 설립중인 재단법인에게 귀속된다. \par
ㄴ. 법인의 불법행위책임이 성립하기 위해서는 대표기관의 행위일 것이 요구되며, 여기서의 대표기관에는 사실상의 대표자도 포함된다. \par
ㄷ. 사단법인 이사의 대표권 제한은 등기되지 않았다고 하더라도 정관에 그 기재가 있는 한, 악의의 제3자에게 대항할 수 있다. \par
ㄹ. 재단법인의 감사는 임의기관이다.
\end{framedtext}
\begin{choices}
  \item ㄱ, ㄴ
  \item ㄱ, ㄷ
  \item ㄴ, ㄷ
  \item ㄴ, ㄹ
  \item ㄷ, ㄹ
\end{choices}
\end{question}
\answer{4}
